\hspace*{5mm}
Darbā ievietoto ilustrāciju - fotogrāfiju, skiču, shēmu, grafiku, diagrammu u. tml. -
apzīmēšanai izmanto vienu un to pašu terminu "attēls". Attēli tiek numurēti nodaļas
ietvaros, un katram no tiem ir jābūt savam nosaukumam. Attēla numuru un
nosaukumu raksta simetriski zem ilustrācijas. Nosaukumu raksta ar maziem burtiem
un lielo sākuma burtu, bez punkta nosaukuma beigās. Trešajā nodaļā ievietota pirmā
pēc kārtas attēla paraksta piemērs:
		\begin{figure}[H]
		  \centering
		  \includegraphics[width=.6\textwidth]{sections/pic/RTU_logotips_rgb_LV.pdf}
		  \caption{RTU logo LV}
		  \label{fig:rtu_lv}
		\end{figure}
\begin{itemize}
    \item Ja attēlam vajadzīgi paskaidrojumi, tos raksta zem nosaukuma (paskaidrojumus var rakstīt, lietojot mazāka izmēra fontu).
    \item Tekstā, kura ilustrēšanai izmantots attēls, attiecīgā vietā jādod atsauce uz to, piemēram: ... \ref{fig:rtu_lv} att. uzskatāmi ilustrē ...
    \item Ja darbā (nodaļā) ir tikai viena ilustrācija, tad to nenumurē, bet raksta tikai tās nosaukumu.	
\end{itemize}

        \begin{figure}[H]
		  \centering
		  \includegraphics[width=.3\textwidth]{sections/pic/RTU_logotips_rgb_EN.pdf}
		  \caption{RTU logo EN}
		  \label{fig:rtu_en}
		\end{figure}

